\SourceLeaf{356}

Schreibt man
{\small\[
(t,t) = a,\quad (t',t') = a',\quad (t'',t'') = a'',\quad (t',t'') = b,\quad (t,t'') = b',\quad (t,t') = b''
\]}
und
\[\begin{aligned}
a'a'' - bb &= A, & aa'' - b'b' &= A', & aa' - b''b'' &= A'', \\
b'b'' - ab &= B, & bb'' - a'b' &= B', & bb' - a''b'' &= B'',
\end{aligned}\]
\[
D = aa'a'' + 2bb'b'' - abb - a'b'b' - a''b''b'',
\]
so gilt
\[
\begin{array}{ll}
 & \text{senkrecht zu}\\
 T=At+B''t'+B't'' & t'\text{ und }t''\\
 T'=B''t+A't'+Bt'' & t\text{ und }t''\\
 T''=B't+Bt'+A''t'' & t\text{ und }t',
\end{array}
\]
und allgemein gilt, wenn
\[
\lambda x + \lambda' x' + \lambda'' x'' = L
\]
die Gleichung einer Ebene ist, dass die Linie
\[
\lambda T + \lambda' T' + \lambda'' T''
\]
senkrecht auf dieser Ebene steht. $L =$ Wert der Form $\begin{pmatrix} A, A', A'' \\ B, B', B'' \end{pmatrix}$,
wenn die Unbestimmten $= \lambda$, $\lambda'$, $\lambda''$ gesetzt werden.

Ferner gilt dann
\[\begin{aligned}
aT + b''T' + b'T'' &= Dt \\
b''T + a'T' + bT'' &= Dt' \\
b'T + bT' + a''T'' &= Dt'',
\end{aligned}\]
und die Linien $t$, $t'$, $t''$ stehen senkrecht auf den Ebenen, deren Gleichungen lauten
\[\begin{aligned}
ax + b''x' + b'x'' &= \text{Const.} \\
b''x + a'x' + bx'' &= \text{Const.} \\
b'x + bx' + a''x'' &= \text{Const.}
\end{aligned}\]

Der doppelte Flächeninhalt des Dreiecks mit den Eckpunkten $mt$, $m't'$, $m''t''$ ist gleich
der Quadratwurzel aus dem Wert der Form
\[
\begin{pmatrix} A, A', A'' \\ B, B', B'' \end{pmatrix} = F(m'm'',\, mm'',\, mm'),
\]
